\documentclass[12pt,letterpaper]{article}
\usepackage{graphicx,textcomp}
\usepackage{natbib}
\usepackage{setspace}
\usepackage{fullpage}
\usepackage{color}
\usepackage[reqno]{amsmath}
\usepackage{amsthm}
\usepackage{fancyvrb}
\usepackage{amssymb,enumerate}
\usepackage[all]{xy}
\usepackage{endnotes}
\usepackage{lscape}
\newtheorem{com}{Comment}
\usepackage{float}
\usepackage{hyperref}
\newtheorem{lem} {Lemma}
\newtheorem{prop}{Proposition}
\newtheorem{thm}{Theorem}
\newtheorem{defn}{Definition}
\newtheorem{cor}{Corollary}
\newtheorem{obs}{Observation}
\usepackage[compact]{titlesec}
\usepackage{dcolumn}
\usepackage{tikz}
\usetikzlibrary{arrows}
\usepackage{multirow}
\usepackage{xcolor}
\newcolumntype{.}{D{.}{.}{-1}}
\newcolumntype{d}[1]{D{.}{.}{#1}}
\definecolor{light-gray}{gray}{0.65}
\usepackage{url}
\usepackage{listings}
\usepackage{color}

\definecolor{codegreen}{rgb}{0,0.6,0}
\definecolor{codegray}{rgb}{0.5,0.5,0.5}
\definecolor{codepurple}{rgb}{0.58,0,0.82}
\definecolor{backcolour}{rgb}{0.95,0.95,0.92}

\lstdefinestyle{mystyle}{
	backgroundcolor=\color{backcolour},   
	commentstyle=\color{codegreen},
	keywordstyle=\color{magenta},
	numberstyle=\tiny\color{codegray},
	stringstyle=\color{codepurple},
	basicstyle=\footnotesize,
	breakatwhitespace=false,         
	breaklines=true,                 
	captionpos=b,                    
	keepspaces=true,                 
	numbers=left,                    
	numbersep=5pt,                  
	showspaces=false,                
	showstringspaces=false,
	showtabs=false,                  
	tabsize=2
}
\lstset{style=mystyle}
\newcommand{\Sref}[1]{Section~\ref{#1}}
\newtheorem{hyp}{Hypothesis}


\title{Problem Set 4}
\date{Due: November 27, 2025}
\author{Applied Stats/Quant Methods 1}


\begin{document}
	\maketitle
	\section*{Instructions}
	\begin{itemize}
		\item Please show your work! You may lose points by simply writing in the answer. If the problem requires you to execute commands in \texttt{R}, please include the code you used to get your answers. Please also include the \texttt{.R} file that contains your code. If you are not sure if work needs to be shown for a particular problem, please ask.
		\item Your homework should be submitted electronically on GitHub.
		\item This problem set is due before 23:59 on Thursday November 27, 2025. No late assignments will be accepted.
	\end{itemize}



	\vspace{.5cm}
\section*{Question 1: Economics}
\vspace{.25cm}
\noindent 	
In this question, use the \texttt{prestige} dataset in the \texttt{car} library. First, run the following commands:

\begin{verbatim}
install.packages(car)
library(car)
data(Prestige)
help(Prestige)
\end{verbatim} 


\noindent We would like to study whether individuals with higher levels of income have more prestigious jobs. Moreover, we would like to study whether professionals have more prestigious jobs than blue and white collar workers.

\newpage
\begin{enumerate}
	
	\item [(a)]
	Create a new variable \texttt{professional} by recoding the variable \texttt{type} so that professionals are coded as $1$, and blue and white collar workers are coded as $0$ (Hint: \texttt{ifelse}).
		\lstinputlisting[language=R, firstline=36, lastline=37]{PS04.R}
		\begin{center}
		\begin{lstlisting}
		This code creates a new binary variable professional where each row is 1 if the worker is a professional and 0 if not. Because some categories are missing through NA I used the prof to identify "prof" in the dataset. 
		\end{lstlisting}
		\end{center}
	\vspace{6cm}
	
	
	\item [(b)]
	Run a linear model with \texttt{prestige} as an outcome and \texttt{income}, \texttt{professional}, and the interaction of the two as predictors (Note: this is a continuous $\times$ dummy interaction.)
		\lstinputlisting[language=R, firstline=38, lastline=39]{PS04.R}
		\begin{center}
		\begin{verbatim}
			Residuals:
			Min      1Q  Median      3Q     Max 
			-32.085  -8.648  -2.498   8.979  31.879 
			
			Coefficients: (2 not defined because of singularities)
			Estimate Std. Error t value Pr(>|t|)    
			(Intercept)         2.760e+01  2.380e+00  11.594  < 2e-16 ***
			income              2.844e-03  2.933e-04   9.694 6.77e-16 ***
			professional               NA         NA      NA       NA    
			income:professional        NA         NA      NA       NA    
			---
			Signif. codes:  0 ‘***’ 0.001 ‘**’ 0.01 ‘*’ 0.05 ‘.’ 0.1 ‘ ’ 1
			
			Residual standard error: 12.22 on 96 degrees of freedom
			(4 observations deleted due to missingness)
			Multiple R-squared:  0.4946,	Adjusted R-squared:  0.4894 
			F-statistic: 93.97 on 1 and 96 DF,  p-value: 6.773e-16
		\end{verbatim}
		\end{center}
	
	\vspace{6cm}
	\item [(c)]
\[
\hat{\text{prestige}} = \beta_0 + \beta_1 \cdot \text{income} + \beta_2 \cdot \text{professional} + \beta_3 \cdot (\text{income} \times \text{professional})
\]


\textit The model predicts prestige using income, professional status, and their interaction. The interaction shows how income’s effect on prestige differs for professionals versus non-professionals, allowing the model to capture how these factors together influence prestige.

\newpage
	\item [(d)]
	Interpret the coefficient for \texttt{income}.


	\textit For non-professionals, each one-unit increase in income is associated with an increase of  0.003 in prestige, all other variables constant. 
	\vspace{10cm}	
	\item [(e)]
	Interpret the coefficient for \texttt{professional}.
	
	
	\textit Professionals score 37.78 points higher in prestige, holding income at zero.  
		\newpage
	\item [(f)]
	What is the effect of a \$1,000 increase in income on prestige score for professional occupations? In other words, we are interested in the marginal effect of income when the variable \texttt{professional} takes the value of $1$. Calculate the change in $\hat{y}$ associated with a \$1,000 increase in income based on your answer for (c).
	
	\lstinputlisting[language=R, firstline=40, lastline=41]{PS04.R}
	
		\textit  A 1,000 increase in income is associated with an 0.84 increase in prestige for people with professional jobs.
		
	
			 
	\vspace{10cm}
	
	
	\item [(g)]
	What is the effect of changing one's occupations from non-professional to professional when her income is \$6,000? We are interested in the marginal effect of professional jobs when the variable \texttt{income} takes the value of $6,000$. Calculate the change in $\hat{y}$ based on your answer for (c).
	
	
	\lstinputlisting[language=R, firstline=42, lastline=43]{PS04.R}
	
	\textit Changing from non-professional to professional at an income of 6,000 increases prestige by 23.82 points. 
	
\end{enumerate}

\newpage

\section*{Question 2: Political Science}
\vspace{.25cm}
\noindent 	Researchers are interested in learning the effect of all of those yard signs on voting preferences.\footnote{Donald P. Green, Jonathan	S. Krasno, Alexander Coppock, Benjamin D. Farrer,	Brandon Lenoir, Joshua N. Zingher. 2016. ``The effects of lawn signs on vote outcomes: Results from four randomized field experiments.'' Electoral Studies 41: 143-150. } Working with a campaign in Fairfax County, Virginia, 131 precincts were randomly divided into a treatment and control group. In 30 precincts, signs were posted around the precinct that read, ``For Sale: Terry McAuliffe. Don't Sellout Virgina on November 5.'' \\

Below is the result of a regression with two variables and a constant.  The dependent variable is the proportion of the vote that went to McAuliff's opponent Ken Cuccinelli. The first variable indicates whether a precinct was randomly assigned to have the sign against McAuliffe posted. The second variable indicates
a precinct that was adjacent to a precinct in the treatment group (since people in those precincts might be exposed to the signs).  \\

\vspace{.5cm}
\begin{table}[!htbp]
	\centering 
	\textbf{Impact of lawn signs on vote share}\\
	\begin{tabular}{@{\extracolsep{5pt}}lccc} 
		\\[-1.8ex] 
		\hline \\[-1.8ex]
		Precinct assigned lawn signs  (n=30)  & 0.042\\
		& (0.016) \\
		Precinct adjacent to lawn signs (n=76) & 0.042 \\
		&  (0.013) \\
		Constant  & 0.302\\
		& (0.011)
		\\
		\hline \\
	\end{tabular}\\
	\footnotesize{\textit{Notes:} $R^2$=0.094, N=131}
\end{table}

\vspace{.5cm}
\begin{enumerate}
	\item [(a)] Use the results from a linear regression to determine whether having these yard signs in a precinct affects vote share (e.g., conduct a hypothesis test with $\alpha = .05$).
	
		\lstinputlisting[language=R, firstline=44, lastline=57]{PS04.R}
		
			\textit Using the outputs from the regression, I took the coefficient for “Precinct assigned lawn signs” (0.042) and its standard error (0.016) to calculate the t-statistic. I then used the t-statistic to determine the p-value and compared it to the significance level, 0.05. Based on this comparison, I rejected the null hypothesis, concluding that there is statistically significant evidence that having lawn signs in a precinct increases vote share. The estimated effect of having lawn signs is 0.042, which corresponds to an approximate 4.2 percent increase in vote share. 
	
	\newpage		
	\item [(b)]  Use the results to determine whether being
	next to precincts with these yard signs affects vote
	share (e.g., conduct a hypothesis test with $\alpha = .05$).
	
	
	\lstinputlisting[language=R, firstline=58, lastline=71]{PS04.R}
	
		\textit Using the outputs from the regression, I took the coefficient for “Precinct assigned lawn signs” (0.042) and its standard error (0.013) to calculate the t-statistic. I then used the t-statistic to determine the p-value and compared it to the significance level, 0.05. Based on this comparison, I rejected the null hypothesis, concluding that there is statistically significant evidence that being in precincts next to lawn signs increases vote share. The estimated effect of being next to lawn signs is 0.042, which corresponds to an approximate 4.2 percent increase in vote share. 
	
	
	\vspace{7cm}
	\item [(c)] Interpret the coefficient for the constant term substantively.
	\vspace{7cm}
	
	\textit The constant term coefficient indicates that when there are no yard signs in a precinct or in adjacent precincts this represents 30.2 percent of the vote share. 
	
	\item [(d)] Evaluate the model fit for this regression.  What does this	tell us about the importance of yard signs versus other factors that are not modeled?
	
		\textit The regression results show that being in a precinct with lawn signs or adjacent to one increases vote share by approximately 4.2 percent, which is statistically significant. However, the baseline vote share of 30.2 percent indicates that most of the variation in vote share is not explained by yard signs. This suggests that, while yard signs have a measurable effect, other factors not included in the model.
	
\end{enumerate}  


\end{document}
